\documentclass[10pt,a4paper]{amsart}
\usepackage{amsmath,amssymb,amsthm,graphicx}
% Retained minimal preamble for the Stacks Project chapter flat.tex (More on Flatness).
% Verbatim excerpt of the chapter preamble: Xy-pic setup (native preamble lines 14-18) and the native
% theorem declarations with the mathematical macros (native preamble lines 154-217). Copied mechanically;
% nothing added, removed or redefined. The wrapper class and the amsmath/amssymb/amsthm/graphicx base
% packages are supplied by the assembler, so the class/externaldocument part of the native preamble is omitted.
% For commutative diagrams we use Xy-pic
\usepackage[all]{xy}

% We use 2cell for 2-commutative diagrams.
\xyoption{2cell}
% Theorem environments.
%
\theoremstyle{plain}
\newtheorem{theorem}[subsection]{Theorem}
\newtheorem{proposition}[subsection]{Proposition}
\newtheorem{lemma}[subsection]{Lemma}

\theoremstyle{definition}
\newtheorem{definition}[subsection]{Definition}
\newtheorem{example}[subsection]{Example}
\newtheorem{exercise}[subsection]{Exercise}
\newtheorem{situation}[subsection]{Situation}

\theoremstyle{remark}
\newtheorem{remark}[subsection]{Remark}
\newtheorem{remarks}[subsection]{Remarks}

\numberwithin{equation}{subsection}

% Macros
%
\def\lim{\mathop{\mathrm{lim}}\nolimits}
\def\colim{\mathop{\mathrm{colim}}\nolimits}
\def\Spec{\mathop{\mathrm{Spec}}}
\def\Hom{\mathop{\mathrm{Hom}}\nolimits}
\def\Ext{\mathop{\mathrm{Ext}}\nolimits}
\def\SheafHom{\mathop{\mathcal{H}\!\mathit{om}}\nolimits}
\def\SheafExt{\mathop{\mathcal{E}\!\mathit{xt}}\nolimits}
\def\Sch{\mathit{Sch}}
\def\Mor{\mathop{\mathrm{Mor}}\nolimits}
\def\Ob{\mathop{\mathrm{Ob}}\nolimits}
\def\Sh{\mathop{\mathit{Sh}}\nolimits}
\def\NL{\mathop{N\!L}\nolimits}
\def\CH{\mathop{\mathrm{CH}}\nolimits}
\def\proetale{{pro\text{-}\acute{e}tale}}
\def\etale{{\acute{e}tale}}
\def\QCoh{\mathit{QCoh}}
\def\Ker{\mathop{\mathrm{Ker}}}
\def\Im{\mathop{\mathrm{Im}}}
\def\Coker{\mathop{\mathrm{Coker}}}
\def\Coim{\mathop{\mathrm{Coim}}}

% Boxtimes
%
\DeclareMathSymbol{\boxtimes}{\mathbin}{AMSa}{"02}

%
% Macros for moduli stacks/spaces
%
\def\QCohstack{\mathcal{QC}\!\mathit{oh}}
\def\Cohstack{\mathcal{C}\!\mathit{oh}}
\def\Spacesstack{\mathcal{S}\!\mathit{paces}}
\def\Quotfunctor{\mathrm{Quot}}
\def\Hilbfunctor{\mathrm{Hilb}}
\def\Curvesstack{\mathcal{C}\!\mathit{urves}}
\def\Polarizedstack{\mathcal{P}\!\mathit{olarized}}
\def\Complexesstack{\mathcal{C}\!\mathit{omplexes}}
% \Pic is the operator that assigns to X its picard group, usage \Pic(X)
% \Picardstack_{X/B} denotes the Picard stack of X over B
% \Picardfunctor_{X/B} denotes the Picard functor of X over B
\def\Pic{\mathop{\mathrm{Pic}}\nolimits}
\def\Picardstack{\mathcal{P}\!\mathit{ic}}
\def\Picardfunctor{\mathrm{Pic}}
\def\Deformationcategory{\mathcal{D}\!\mathit{ef}}

% --- appended 2026-10-12 (Chief): the project own definitions of the reference/slogan/history
% display environments, copied mechanically from the shared Stacks preamble lines 8-12
% (\usepackage{verbatim} plus the three \newenvironment lines, which the chapter preamble keeps
% verbatim from the project preamble). Nothing added, removed or redefined: these environments
% carry no mathematical content and the project itself typesets them as comments. Without these
% three lines a counted statement whose own text uses them fails with Environment slogan undefined.
\usepackage{verbatim}
\newenvironment{reference}{\comment}{\endcomment}
\newenvironment{slogan}{\comment}{\endcomment}
\newenvironment{history}{\comment}{\endcomment}
\usepackage{cleveref}
\crefname{originallemma}{original lemma}{original lemmas}
\Crefname{originallemma}{Original lemma}{Original lemmas}
\makeatletter
\newlabel{morphisms-lemma-monomorphism-loc-finite-type-loc-quasi-finite}{{0CT8}{source}{}{sourcepointer-1}{}}
\newlabel{morphisms-lemma-monomorphism-loc-finite-type-loc-quasi-finite@cref}{{[originallemma][0][]0CT8}{[1][0][]source}{}{}{}}
\makeatother
\crefname{originallemma}{original lemma}{original lemmas}
\Crefname{originallemma}{Original lemma}{Original lemmas}
\makeatletter
\newlabel{more-morphisms-lemma-quasi-finite-separated-pass-through-finite}{{05K0}{source}{}{sourcepointer-2}{}}
\newlabel{more-morphisms-lemma-quasi-finite-separated-pass-through-finite@cref}{{[originallemma][0][]05K0}{[1][0][]source}{}{}{}}
\makeatother
\crefname{originallemma}{original lemma}{original lemmas}
\Crefname{originallemma}{Original lemma}{Original lemmas}
\makeatletter
\newlabel{more-algebra-lemma-limit-henselian}{{0CT7}{source}{}{sourcepointer-3}{}}
\newlabel{more-algebra-lemma-limit-henselian@cref}{{[originallemma][0][]0CT7}{[1][0][]source}{}{}{}}
\makeatother
\title{A locally finite type monomorphism into the spectrum of a limit ring that receives all finite stages is an isomorphism}
\author{The Stacks Project Authors}
\date{stacks-project:flat.tex:local-snapshot-20261011; GNU FDL 1.2 or later}
\begin{document}
\maketitle
\noindent\textit{Editorial note.} All delivered mathematical text is the author's own, copied from the pinned Stacks Project chapter file flat.tex; no statement or proof step was rewritten, reordered or removed.\par\medskip
\noindent\textit{Editorial note.} The preamble is the chapter's own retained minimal preamble (sources/stacks-flat/retained-preamble.tex), a verbatim excerpt of the native Xy-pic setup and of the native theorem declarations and mathematical macros, followed by the project own definitions of the reference, slogan and history display environments (verbatim package plus three newenvironment lines) that the retained file carries at its end and that the project itself typesets as comments. No macro was invented or redefined, and the wrapper class and the amsmath/amssymb/amsthm/graphicx base packages are supplied by the assembler.\par\medskip
\noindent\textit{Editorial note.} The counted statement and proof contain no cross reference, no citation and no included file, so the delivered body is self contained.\par\medskip
\noindent\textit{Editorial note.} Canonical location of the counted proposition: Stacks tag 0CTJ (sources/stacks-flat/pinned-original-source.tex line 8274).\par\medskip
\noindent\textit{Allowed standard prerequisites (editorial).}\begin{itemize}
\item A monomorphism locally of finite type is separated and quasi-finite. Morphisms, Lemma labelled morphisms-lemma-monomorphism-loc-finite-type-loc-quasi-finite (Stacks tag 0CT8), invoked at proof lines 8292-8294. Supplies that the quasi-compact T -> S is separated and quasi-finite, which makes Zariski's Main Theorem applicable.
\item Zariski's Main Theorem in the quasi-finite separated form. More on Morphisms, Lemma labelled more-morphisms-lemma-quasi-finite-separated-pass-through-finite (Stacks tag 05K0), invoked at proof lines 8295-8298. Supplies the factorization T -> W -> S with W finite over S and T -> W an open immersion, which converts the problem into a statement about the ring map A -> B.
\item The pair (A, Ker(A -> A\_n)) is henselian for a limit with surjective transition maps and locally nilpotent kernels. More on Algebra, Lemma labelled more-algebra-lemma-limit-henselian (Stacks tag 0CT7), invoked at proof lines 8313-8314. Supplies that every closed point of S lies in the image of S\_n -> S, which is what forces an open containing that image to be everything.
\end{itemize}\medskip
\ifdefined\hypertarget\hypertarget{sourcepointer-1}{}\fi
\noindent\textit{Source reference.} Original lemma 0CT8: sources/stacks-chow/tags.txt, source lines 16058--16058. Cross-chapter standard prerequisite located in the project tag list staged in the local snapshot; invoked at proof line 8294 of the counted proof. It is retained only as a navigation pointer; every load-bearing step is inside the delivered span.\par\medskip
\ifdefined\hypertarget\hypertarget{sourcepointer-2}{}\fi
\noindent\textit{Source reference.} Original lemma 05K0: sources/stacks-chow/tags.txt, source lines 7184--7184. Cross-chapter standard prerequisite located in the project tag list staged in the local snapshot; invoked at proof line 8297 of the counted proof. It is retained only as a navigation pointer; every load-bearing step is inside the delivered span.\par\medskip
\ifdefined\hypertarget\hypertarget{sourcepointer-3}{}\fi
\noindent\textit{Source reference.} Original lemma 0CT7: sources/stacks-chow/tags.txt, source lines 16057--16057. Cross-chapter standard prerequisite located in the project tag list staged in the local snapshot; invoked at proof line 8314 of the counted proof. It is retained only as a navigation pointer; every load-bearing step is inside the delivered span.\par\medskip
\begin{lemma}
\label{lemma-monomorphism-isomorphism}
Let $A = \lim A_n$ be a limit of a system of rings
whose transition maps are surjective and with locally nilpotent
kernels. Let $S = \Spec(A)$. Let $T \to S$ be a monomorphism
which is locally of finite type. If $\Spec(A_n) \to S$
factors through $T$ for all $n$, then $T = S$.
\end{lemma}

\begin{proof}
Set $S_n = \Spec(A_n)$. Let $T_0 \subset T$ be the common
image of the factorizations $S_n \to T$. Then $T_0$ is quasi-compact.
Let $T' \subset T$ be a quasi-compact open containing $T_0$.
Then $S_n \to T$ factors through $T'$.
If we can show that $T' = S$, then $T' = T = S$.
Hence we may assume $T$ is quasi-compact.

\medskip\noindent
Assume $T$ is quasi-compact.
In this case $T \to S$ is separated and quasi-finite
(Morphisms, Lemma
\ref{morphisms-lemma-monomorphism-loc-finite-type-loc-quasi-finite}).
Using Zariski's Main Theorem
(in the form of More on Morphisms, Lemma
\ref{more-morphisms-lemma-quasi-finite-separated-pass-through-finite})
we choose a factorization $T \to W \to S$ with $W \to S$ finite
and $T \to W$ an open immersion. Write $W = \Spec(B)$.
The (unique) factorizations $S_n \to T$ may be viewed
as morphisms into $W$ and we obtain
$$
A \longrightarrow B \longrightarrow \lim A_n = A
$$
Consider the morphism $h : S = \Spec(A) \to \Spec(B) = W$ coming
from the arrow on the right. Then
$$
T \times_{W, h} S
$$
is an open subscheme of $S$ containing the image of $S_n \to S$ for all $n$.
To finish the proof it suffices to show that any open $U \subset S$
containing the image of $S_n \to S$ for some $n \geq 1$ is equal to $S$.
This is true because $(A, \Ker(A \to A_n))$ is a henselian pair
(More on Algebra, Lemma \ref{more-algebra-lemma-limit-henselian})
and hence every closed point of $S$ is contained in the image of $S_n \to S$.
\end{proof}

\end{document}
